\documentclass[12pt, a4paper]{article}
\usepackage[utf8]{inputenc}
\usepackage{amsmath, amssymb, amsthm, amsfonts}
\usepackage{CJKutf8}
\usepackage[margin=1in]{geometry}
\usepackage{tikz}
\newtheorem{lemma}{Lemma}
\newcommand{\g}{\mathfrak{g}}
\newcommand{\h}{\mathfrak{h}}
\newcommand{\R}{\mathbb{R}}
% 重新定义 solution 环境，保留标识并在末尾留出 8cm 空白
\newenvironment{solution}
  {\paragraph{\textit{Solution:}}\par\medskip}
  {\hfill$\blacksquare$\vspace{0.1cm}}

\title{Lie Theory:Assignment 5}
\author{
    \begin{CJK*}{UTF8}{gbsn}
    钟皓宇 (12533556)
    \end{CJK*}
}
\date{\today}

\begin{document}

\maketitle




\section*{Problem 1}
Find Weyl groups of $\mathrm{A}_2, \mathrm{B}_2, \mathrm{G}_2$.
\begin{solution}
  Suppose the root system we want to discuss has simple roots $\beta_1, \beta_2$ with angle $\theta$, and let its Weyl group be denoted by $W$. Let $\theta = \pi - \pi/m$.
  Let $s_1, s_2$ denote the reflections corresponding to $\beta_1, \beta_2$ respectively. 
  According to our textbook, the Weyl group can be generated by the reflections of the simple roots; therefore, we just need to clarify the relations between $s_1$ and $s_2$.
  
  For $A_2$, since $\theta = 2\pi/3$ and $m = 3$, we then have:$$s_1(\beta_1) = -\beta_1 \quad s_1(\beta_2) = \beta_1 + \beta_2 \quad s_2(\beta_1) = \beta_1 + \beta_2 \quad s_2(\beta_2) = -\beta_2.$$
  Then we obtain that $(s_1s_2)^3 = \text{id}$. Therefore, the Weyl group of $A_2$ is$$W = \langle s_1, s_2 \mid s_1^2 = s_2^2 = (s_1s_2)^3 = \text{id} \rangle$$
  which is exactly the Dihedral group $D_3$.

  Similarly, if we suppose that $\beta_1$ is the long root of $\mathrm{B}_2$ then we have 
  $$s_1(\beta_1) = -\beta_1 \quad s_1(\beta_2) = \beta_1 + \beta_2 \quad s_2(\beta_1) = \beta_1 + 2\beta_2 \quad s_2(\beta_2) = -\beta_2$$
  and thus $(s_1s_2)^{4}=\text{id}$. Then the Weyl group of $\mathrm{B}_2$ is $W = D_4$.

  For $\mathrm{G}_2$, suppose that $\beta_1$ is the long root again. The following computation 
  $$s_1(\beta_1) = -\beta_1 \quad s_1(\beta_2) = \beta_1 + \beta_2 \quad s_2(\beta_1) = \beta_1 + 3\beta_2 \quad s_2(\beta_2) = -\beta_2$$
  implies that $(s_1s_2)^{6}=\text{id}$ and thus the Weyl group of $\mathrm{G}_2$ is $W = D_6$.
\end{solution}

\section*{Problem 2}
Let $\Delta \subset \mathbb{R}^n$ be a root system and let $B = \{\beta_1, \dots, \beta_n\}$ be a base of $\Delta$. The Cartan matrix is defined as
$$
\left( \frac{2(\beta_i|\beta_j)}{(\beta_j|\beta_j)} \right)_{1 \le i, j \le n}.
$$
\begin{enumerate}
  \item Prove that the Cartan matrix of a root system is symmetrizable positive definite.
  \item Find all Cartan matrices of rank 3.
  \item Show that a Cartan matrix determines the root system uniquely up to isomorphism.
\end{enumerate}


\begin{solution}
  (1) We let $A=(\beta_i|\beta_j)_{1 \le i, j \le n}$ and $x=\sum_{i}a_i \beta_i$ be an arbitrary nonzero vector in $\R^n$. Then the following result 
  \begin{align*}
    x^T A x = \left(\sum_i a_i\beta_i\mid \sum_i a_i\beta_i\right) > 0
  \end{align*}
  implies that $A$ is positive definite. 

  (2) Suppose the Dynkin diagram is disconnected. Equivalently, we say the root system is reducible. Then the there are two possible cases of Dynkin diagram:
  \begin{enumerate}
    \item Totally disconnected: $\Delta = \Delta_1\sqcup \Delta_2 \sqcup \Delta_3$ where $\Delta_i$ is the 1-dimensional root system.
    \item Partially disconnected: $\Delta= \Delta_1 \sqcup \Delta_2$ where $\Delta_1$ is the 1-dimensional root system and $\Delta_2$ be the irreducible 2-dimensional root system.
  \end{enumerate}
  In the totally disconnected case, the Cartan matrix is
\[
A=2I_3=
\begin{pmatrix}
2&0&0\\
0&2&0\\
0&0&2
\end{pmatrix}.
\]
In the case with two connected components, after reordering the simple roots, the Cartan matrix has block diagonal form
\[
A=(2)\oplus A_{\Delta_2}.
\]
Since the irreducible rank-two root systems are of type \(A_2\), \(B_2=C_2\), and \(G_2\), the possible Cartan matrices up to simultaneous permutation of the simple roots are
\[
\begin{pmatrix}
2&0&0\\
0&2&-1\\
0&-1&2
\end{pmatrix},
\qquad
\begin{pmatrix}
2&0&0\\
0&2&-1\\
0&-2&2
\end{pmatrix},
\qquad
\begin{pmatrix}
2&0&0\\
0&2&-1\\
0&-3&2
\end{pmatrix}.
\]

Now we suppose that $\Delta$ is an irreducible rank 3 root system. Then we have two possible cases 
\begin{enumerate}
  \item Any two roots are not orthogonal.
  \item Exists two roots are orthogonal.
\end{enumerate} 
For case 1, we can suppose that the Cartan matrix is 
\begin{align*}
  A=
\begin{pmatrix}
2&-a_{12}&-a_{13}\\
-a_{21}&2&-a_{23}\\
-a_{31}&-a_{32}&2
\end{pmatrix}
\end{align*}
where $a_{ij}>0$. The determinant of $A$ is 
\begin{align*}
  \det(A) = 8
-2(a_{12}a_{21}+a_{13}a_{31}+a_{23}a_{32})
-(a_{12}a_{23}a_{31}+a_{13}a_{32}a_{21})\le 0.
\end{align*}
However, (1) implies that $\det(A)>0$ which is a contradiction. Therefore, the case 1 is impossible.
Consider the case 2. By permutation of the simple roots we can suppose the Cartan matrix be the form of 
\begin{align*}
  A=
\begin{pmatrix}
2&-a_{12}&0\\
-a_{21}&2&-a_{23}\\
0&-a_{32}&2
\end{pmatrix}
\end{align*}
and with determinant 
\begin{align*}
  \det(A) = 8
-2(a_{12}a_{21}+a_{23}a_{32})
\end{align*}
Since we reuqire $\det(A)>0$, it forces 
\begin{align*}
  0<a_{12}a_{21}+a_{23}a_{32}<4.
\end{align*}
We can suppose that $a_{12}a_{21}\le a_{23}a_{32}$ by permutation of simple roots. For any pair of distinct simple roots, we have
\[
a_{ij}a_{ji}\in\{0,1,2,3\}.
\]
Since we suppose the root system $\Delta$ is irreducible, equivalently, both $a_{12}a_{21}$ and $a_{23}a_{32}$ are not equal to zero.
It implies that 
\begin{align*}
  (a_{12}a_{21},a_{23}a_{32}) = (1,1)\text{ or }(1,2).
\end{align*}
which the Cartan matrix of this two case corresponds to $\mathrm{A}_2$ and $\{\mathrm{B}_2,\mathrm{C}_2\}$. Therefore, the Cartan matrix of irreducible rank 3 root system has two possible form 
\begin{align*}
  \begin{pmatrix}
2&-1&0\\
-1&2&-1\\
0&-1&2
\end{pmatrix}, \qquad 
\begin{pmatrix}
2&-1&0\\
-1&2&-2\\
0&-1&2
\end{pmatrix}, \qquad 
\begin{pmatrix}
2&-1&0\\
-1&2&-1\\
0&-2&2
\end{pmatrix}.
\end{align*}

(3) Let \(\Delta\subset V\) and \(\Delta'\subset V'\) be two root systems with bases
\[
B=\{\beta_1,\dots,\beta_n\},
\qquad
B'=\{\beta'_1,\dots,\beta'_n\}.
\]
Assume that they have the same Cartan matrix
\[
A=(a_{ij}).
\]
Define a linear map
\[
T:V\longrightarrow V'
\]
by
\[
T(\beta_i)=\beta'_i.
\]
Since the simple roots form a basis of \(V\), this determines \(T\) uniquely.

Let \(s_j\) be the reflection corresponding to \(\beta_j\), and let \(s'_j\) be the
reflection corresponding to \(\beta'_j\). By the definition of the Cartan matrix,
we have
\[
s_j(\beta_i)=\beta_i-a_{ij}\beta_j.
\]
Similarly,
\[
s'_j(\beta'_i)=\beta'_i-a_{ij}\beta'_j.
\]
Hence
\[
T(s_j(\beta_i))
=
T(\beta_i-a_{ij}\beta_j)
=
\beta'_i-a_{ij}\beta'_j
=
s'_j(\beta'_i)
=
s'_j(T(\beta_i)).
\]
Since this holds for every basis vector \(\beta_i\), we get
\[
T\circ s_j=s'_j\circ T
\]
for every \(j\).

Let \(W\) and \(W'\) be the Weyl groups of \(\Delta\) and \(\Delta'\). The above
identity implies that \(T\) identifies the action of \(W\) with the action of
\(W'\). Therefore, for every word \(w=s_{i_1}\cdots s_{i_k}\in W\), we have
\[
T(w\beta_i)=w'\beta'_i,
\]
where \(w'=s'_{i_1}\cdots s'_{i_k}\in W'\).

Now use the standard fact that every root is obtained from a simple root by the
action of the Weyl group:
\[
\Delta=W\cdot B,
\qquad
\Delta'=W'\cdot B'.
\]
It follows that
\[
T(\Delta)=\Delta'.
\]
Thus \(T\) is an isomorphism of root systems. Hence the Cartan matrix determines
the root system uniquely up to isomorphism.

\end{solution}



\end{document}